% Options for packages loaded elsewhere
% Options for packages loaded elsewhere
\PassOptionsToPackage{unicode}{hyperref}
\PassOptionsToPackage{hyphens}{url}
\PassOptionsToPackage{dvipsnames,svgnames,x11names}{xcolor}
%
\documentclass[
  letterpaper,
  DIV=11,
  numbers=noendperiod]{scrartcl}
\usepackage{xcolor}
\usepackage{amsmath,amssymb}
\setcounter{secnumdepth}{5}
\usepackage{iftex}
\ifPDFTeX
  \usepackage[T1]{fontenc}
  \usepackage[utf8]{inputenc}
  \usepackage{textcomp} % provide euro and other symbols
\else % if luatex or xetex
  \usepackage{unicode-math} % this also loads fontspec
  \defaultfontfeatures{Scale=MatchLowercase}
  \defaultfontfeatures[\rmfamily]{Ligatures=TeX,Scale=1}
\fi
\usepackage{lmodern}
\ifPDFTeX\else
  % xetex/luatex font selection
\fi
% Use upquote if available, for straight quotes in verbatim environments
\IfFileExists{upquote.sty}{\usepackage{upquote}}{}
\IfFileExists{microtype.sty}{% use microtype if available
  \usepackage[]{microtype}
  \UseMicrotypeSet[protrusion]{basicmath} % disable protrusion for tt fonts
}{}
\makeatletter
\@ifundefined{KOMAClassName}{% if non-KOMA class
  \IfFileExists{parskip.sty}{%
    \usepackage{parskip}
  }{% else
    \setlength{\parindent}{0pt}
    \setlength{\parskip}{6pt plus 2pt minus 1pt}}
}{% if KOMA class
  \KOMAoptions{parskip=half}}
\makeatother
% Make \paragraph and \subparagraph free-standing
\makeatletter
\ifx\paragraph\undefined\else
  \let\oldparagraph\paragraph
  \renewcommand{\paragraph}{
    \@ifstar
      \xxxParagraphStar
      \xxxParagraphNoStar
  }
  \newcommand{\xxxParagraphStar}[1]{\oldparagraph*{#1}\mbox{}}
  \newcommand{\xxxParagraphNoStar}[1]{\oldparagraph{#1}\mbox{}}
\fi
\ifx\subparagraph\undefined\else
  \let\oldsubparagraph\subparagraph
  \renewcommand{\subparagraph}{
    \@ifstar
      \xxxSubParagraphStar
      \xxxSubParagraphNoStar
  }
  \newcommand{\xxxSubParagraphStar}[1]{\oldsubparagraph*{#1}\mbox{}}
  \newcommand{\xxxSubParagraphNoStar}[1]{\oldsubparagraph{#1}\mbox{}}
\fi
\makeatother


\usepackage{longtable,booktabs,array}
\usepackage{calc} % for calculating minipage widths
% Correct order of tables after \paragraph or \subparagraph
\usepackage{etoolbox}
\makeatletter
\patchcmd\longtable{\par}{\if@noskipsec\mbox{}\fi\par}{}{}
\makeatother
% Allow footnotes in longtable head/foot
\IfFileExists{footnotehyper.sty}{\usepackage{footnotehyper}}{\usepackage{footnote}}
\makesavenoteenv{longtable}
\usepackage{graphicx}
\makeatletter
\newsavebox\pandoc@box
\newcommand*\pandocbounded[1]{% scales image to fit in text height/width
  \sbox\pandoc@box{#1}%
  \Gscale@div\@tempa{\textheight}{\dimexpr\ht\pandoc@box+\dp\pandoc@box\relax}%
  \Gscale@div\@tempb{\linewidth}{\wd\pandoc@box}%
  \ifdim\@tempb\p@<\@tempa\p@\let\@tempa\@tempb\fi% select the smaller of both
  \ifdim\@tempa\p@<\p@\scalebox{\@tempa}{\usebox\pandoc@box}%
  \else\usebox{\pandoc@box}%
  \fi%
}
% Set default figure placement to htbp
\def\fps@figure{htbp}
\makeatother





\setlength{\emergencystretch}{3em} % prevent overfull lines

\providecommand{\tightlist}{%
  \setlength{\itemsep}{0pt}\setlength{\parskip}{0pt}}



 


\KOMAoption{captions}{tableheading}
\makeatletter
\@ifpackageloaded{caption}{}{\usepackage{caption}}
\AtBeginDocument{%
\ifdefined\contentsname
  \renewcommand*\contentsname{Table of contents}
\else
  \newcommand\contentsname{Table of contents}
\fi
\ifdefined\listfigurename
  \renewcommand*\listfigurename{List of Figures}
\else
  \newcommand\listfigurename{List of Figures}
\fi
\ifdefined\listtablename
  \renewcommand*\listtablename{List of Tables}
\else
  \newcommand\listtablename{List of Tables}
\fi
\ifdefined\figurename
  \renewcommand*\figurename{Figure}
\else
  \newcommand\figurename{Figure}
\fi
\ifdefined\tablename
  \renewcommand*\tablename{Table}
\else
  \newcommand\tablename{Table}
\fi
}
\@ifpackageloaded{float}{}{\usepackage{float}}
\floatstyle{ruled}
\@ifundefined{c@chapter}{\newfloat{codelisting}{h}{lop}}{\newfloat{codelisting}{h}{lop}[chapter]}
\floatname{codelisting}{Listing}
\newcommand*\listoflistings{\listof{codelisting}{List of Listings}}
\usepackage{amsthm}
\theoremstyle{remark}
\AtBeginDocument{\renewcommand*{\proofname}{Proof}}
\newtheorem*{remark}{Remark}
\newtheorem*{solution}{Solution}
\newtheorem{refremark}{Remark}[section]
\newtheorem{refsolution}{Solution}[section]
\makeatother
\makeatletter
\makeatother
\makeatletter
\@ifpackageloaded{caption}{}{\usepackage{caption}}
\@ifpackageloaded{subcaption}{}{\usepackage{subcaption}}
\makeatother
\usepackage{bookmark}
\IfFileExists{xurl.sty}{\usepackage{xurl}}{} % add URL line breaks if available
\urlstyle{same}
\hypersetup{
  pdftitle={Pearson Correlation for Ordinal Variables: Consolidated Analysis},
  colorlinks=true,
  linkcolor={blue},
  filecolor={Maroon},
  citecolor={Blue},
  urlcolor={Blue},
  pdfcreator={LaTeX via pandoc}}


\title{Pearson Correlation for Ordinal Variables: Consolidated Analysis}
\author{}
\date{}
\begin{document}
\maketitle

\renewcommand*\contentsname{Table of contents}
{
\hypersetup{linkcolor=}
\setcounter{tocdepth}{3}
\tableofcontents
}

\section{Extremal Pearson Correlation for Ordinal
Variables}\label{extremal-pearson-correlation-for-ordinal-variables}

\subsection{1. Problem Setting and
Notation}\label{problem-setting-and-notation}

We consider two ordinal variables:

\begin{itemize}
\tightlist
\item
  \(X\) with \(K_X\) ordered categories \(0, 1, \dots, K_X-1\)
\item
  \(Y\) with \(K_Y\) ordered categories \(0, 1, \dots, K_Y-1\)
\end{itemize}

\textbf{Marginals} given by:

\[
p_X(i) = P(X = i), \quad p_Y(j) = P(Y = j)
\] or counts: \[
n_{X=i}, \quad n_{Y=j}, \quad N = \sum_i n_{X=i} = \sum_j n_{Y=j}
\]

\textbf{CDFs}: \[
F_X(i) = \sum_{k \le i} p_X(k), \quad F_Y(j) = \sum_{k \le j} p_Y(k)
\]

\textbf{Pearson correlation}: \[
r_{XY} = \frac{\mathrm{Cov}(X,Y)}{\sigma_X \sigma_Y}, \quad
E[XY] = \sum_{i=0}^{K_X-1} \sum_{j=0}^{K_Y-1} x_i y_j \, P(X=x_i, Y=y_j)
\] For Spearman's correlation, replace \(X,Y\) by their ranks
\(R_X,R_Y\).

\emph{(Sources: Proofs-OrdinalVariables.md; Ordinal Variables and
Correlation.md; BFbounds.md; Math Proof GPT -- Whitt 1976 Paper.md)}

\begin{center}\rule{0.5\linewidth}{0.5pt}\end{center}

\subsection{2. Fréchet--Hoeffding / Boole--Fréchet
Bounds}\label{fruxe9chethoeffding-boolefruxe9chet-bounds}

Any joint CDF \(F_{XY}\) with marginals \(F_X, F_Y\) satisfies:

\[
\max\{F_X(i) + F_Y(j) - 1, 0\} \ \le \ F_{XY}(i,j) \ \le \ \min\{F_X(i), F_Y(j)\}
\]\\
\emph{(BFbounds.md)}

\begin{itemize}
\tightlist
\item
  \textbf{Upper bound (comonotonic)}: \[
  F^{\text{upper}}_{XY}(i,j) = \min\{F_X(i), F_Y(j)\}
  \]
\item
  \textbf{Lower bound (countermonotonic)}: \[
  F^{\text{lower}}_{XY}(i,j) = \max\{F_X(i) + F_Y(j) - 1, 0\}
  \]
\end{itemize}

\begin{center}\rule{0.5\linewidth}{0.5pt}\end{center}

\subsection{\texorpdfstring{3. Maximizing and Minimizing
\(r\)}{3. Maximizing and Minimizing r}}\label{maximizing-and-minimizing-r}

\paragraph{\texorpdfstring{3.1 Max \(r\) --- Comonotonic
Coupling}{3.1 Max r --- Comonotonic Coupling}}\label{max-r-comonotonic-coupling}

\textbf{Claim:} The maximum correlation \(r_{\max}\) is achieved by the
comonotonic (FH upper) arrangement.

\textbf{Proofs / Arguments:} - \textbf{FH upper pmf construction}: In
discrete case, differences of CDF steps allocate mass along the
diagonal, pairing high \(X\) with high \(Y\).\\
\emph{(BFbounds.md; Proofs-OrdinalVariables.md)} - \textbf{Rearrangement
inequality (ranks)}: For ranks \(R_X, R_Y\), \(\sum R_X R_Y\) is
maximized when sequences are ordered the same way.\\
\emph{(Ordinal Variables and Correlation.md)} - \textbf{Quantile
coupling}: \(U \sim \mathrm{Unif}(0,1)\),
\(X = F_X^{-1}(U), Y = F_Y^{-1}(U)\).\\
\emph{(Math Proof GPT -- Whitt 1976 Paper.md)}

\begin{center}\rule{0.5\linewidth}{0.5pt}\end{center}

\paragraph{\texorpdfstring{3.2 Min \(r\) --- Countermonotonic
Coupling}{3.2 Min r --- Countermonotonic Coupling}}\label{min-r-countermonotonic-coupling}

\textbf{Claim:} The minimal correlation \(r_{\min}\) corresponds
conceptually to the countermonotonic (FH lower) arrangement.

\textbf{Proofs / Arguments:} - \textbf{FH lower pmf construction}:
Differences of the lower-bound CDF.\\
\emph{(BFbounds.md)} - \textbf{Opposite-order pairing}: Sort \(X\)
descending, \(Y\) ascending; achieves best feasible negative
correlation.\\
\emph{(Proofs-OrdinalVariables.md; Ordinal Variables and
Correlation.md)} - \textbf{Quantile coupling}:
\(X = F_X^{-1}(U), Y = F_Y^{-1}(1-U)\) in ideal continuous case.\\
\emph{(Math Proof GPT -- Whitt 1976 Paper.md)}

\textbf{Note:} In discrete settings, FH lower may not be attainable; see
Section 5.

\begin{center}\rule{0.5\linewidth}{0.5pt}\end{center}

\subsection{4. Binary (2×2) Cases and
Examples}\label{binary-22-cases-and-examples}

\emph{(From Cheemera v1.1 -- Pearson Correlation Analysis
Implications.md)}

\paragraph{4.1 Balanced marginals (70--30,
70--30)}\label{balanced-marginals-7030-7030}

\begin{itemize}
\tightlist
\item
  Perfect diagonal fill yields \(r_{\max} = 1\).
\end{itemize}

\paragraph{4.2 Skew vs balanced (90--10,
50--50)}\label{skew-vs-balanced-9010-5050}

\begin{itemize}
\tightlist
\item
  Perfect alignment impossible; \(r_{\max} < 1\) due to capacity limits.
\end{itemize}

\begin{center}\rule{0.5\linewidth}{0.5pt}\end{center}

\subsection{5. Attainability and Symmetry
(Updated)}\label{attainability-and-symmetry-updated}

\subsubsection{5.1 Maximal correlation}\label{maximal-correlation}

\begin{itemize}
\tightlist
\item
  \textbf{Result:} In the discrete ordinal setting, \(r_{\max}\) is
  \textbf{always} attained by the FH upper (comonotonic) construction.\\
  \emph{(BFbounds.md; Proofs-OrdinalVariables.md; Ordinal Variables and
  Correlation.md)}
\end{itemize}

\subsubsection{5.2 Minimal correlation}\label{minimal-correlation}

\begin{itemize}
\tightlist
\item
  \textbf{Result:} The FH lower (countermonotonic) construction may be
  \textbf{unattainable} in discrete settings; \(r_{\min}\) is then the
  best feasible opposite-order arrangement.\\
  \emph{(BFbounds.md; Proofs-OrdinalVariables.md; Ordinal Variables and
  Correlation.md)}
\end{itemize}

\subsubsection{5.3 Magnitude
relationships}\label{magnitude-relationships}

\begin{itemize}
\tightlist
\item
  \textbf{Earlier view:} Often \(|r_{\min}| < r_{\max}\) for discrete
  asymmetrical marginals.\\
  \emph{(BFbounds.md; Ordinal Variables and Correlation.md)}
\item
  \textbf{New evidence (BES data)}: Possible to have
  \(|r_{\min}| > r_{\max}|\).\\
  Example: \(X\) with 4 categories, \(Y\) with 5 categories, skewed
  marginals;\\
  \(r_{\max} = 0.85027\), \(r_{\min} = -0.89916\).\\
  \emph{(Cees -- r\_min MAGNITUDE Asymmetry in Pearson Correlation.md)}
\end{itemize}

\textbf{Reason:} Pearson correlation uses centered values; skewness can
make negative deviations in antitone pairing larger in magnitude than
positive deviations in comonotone pairing.

\subsubsection{Lemma (signs of extreme correlations under monotone
scoring)}\label{lemma-signs-of-extreme-correlations-under-monotone-scoring}

\phantomsection\label{prop-max-product}
\begin{lemma}

Let $X$ and $Y$ be ordinal variables. Assign numeric scores to categories via **non-decreasing** maps $s_X, s_Y$ that respect the category order (scores may be affine or nonlinear; ties allowed). Define two extreme couplings:

* **Comonotone (upper)**: pair larger scored values of $X$ with larger scored values of $Y$.
    
* **Antitone (lower)**: pair larger scored values of $X$ with smaller scored values of $Y$.
    

Let $r_{\max}$ and $r_{\min}$ be the Pearson correlations computed from these couplings, either (i) in the **finite-sample/permutation** sense (sort scored observations in the same vs. opposite order), or (ii) in the **Fréchet–Hoeffding (FH)/population** sense (pair quantiles $Q_X(u)$ with $Q_Y(u)$ vs. $Q_Y(1-u)$).

Then:

1. $r_{\max}\ge 0$ and $r_{\min}\le 0$.
    
2. Moreover, $r_{\max}>0$ and $r_{\min}<0$ **iff** both variables have **positive variance** under the chosen scoring (i.e., neither scored variable is almost surely constant).
    

_(If you reverse the scoring of one variable (decreasing map), swap “upper” and “lower”. For Spearman’s $$\rho$$, take the scores to be midranks—monotone by construction—and the same conclusions hold.)_
\end{lemma}

\begin{proof}
Proof

We give parallel proofs in the finite-sample and population settings.

\begin{enumerate}
\def\labelenumi{\Alph{enumi})}
\tightlist
\item
  Finite-sample (permutation) setting
\end{enumerate}

Let \(x_{(1)}\ge\cdots\ge x_{(n)}\) and \(y_{(1)}\ge\cdots\ge y_{(n)}\)
be the \textbf{scored} observations of \(X\) and \(Y\) sorted in
nonincreasing order (monotone scoring ensures sorting by category equals
sorting by score). Center them:

\[x'_i := x_{(i)}-\bar x,\qquad y'_i := y_{(i)}-\bar y,\quad  
\bar x := \tfrac1n\sum_i x_{(i)},\;\bar y := \tfrac1n\sum_i y_{(i)}.\]

Define the two extreme centered sums

\[S_{\max} := \sum_{i=1}^n x'_i\,y'_i,\qquad  
S_{\min} := \sum_{i=1}^n x'_i\,y'_{\,n+1-i}.\]

\begin{itemize}
\item
  \textbf{Signs.}\\
  Chebyshev's sum inequality (rearrangement for similarly/ oppositely
  ordered sequences) gives

  \[\frac1n\sum_i x'_i y'_i \;\ge\; \Big(\tfrac1n\sum_i x'_i\Big)\Big(\tfrac1n\sum_i y'_i\Big)=0,  
    \quad\Rightarrow\; S_{\max}\ge 0,\]

  and, because \(y'_{\,n+1-i}\) is \textbf{nondecreasing} in \(i\),

  \[\frac1n\sum_i x'_i y'_{\,n+1-i} \;\le\; \Big(\tfrac1n\sum_i x'_i\Big)\Big(\tfrac1n\sum_i y'_{\,n+1-i}\Big)=0,  
    \quad\Rightarrow\; S_{\min}\le 0.\]

  Dividing by \(\sigma_X\sigma_Y>0\) (the standard deviations of the
  scored variables) yields \(r_{\max}\ge 0\) and \(r_{\min}\le 0\).
\item
  \textbf{Strictness \((\Leftrightarrow\) positive variances).}\\
  Suppose \(\sigma_X>0\) and \(\sigma_Y>0\) (i.e., neither \(\{x'_i\}\)
  nor \(\{y'_i\}\) is identically zero). The rearrangement inequality
  says \(S_{\max}\) is the \textbf{maximum} and \(S_{\min}\) the
  \textbf{minimum} of \(\sum_i x'_i y'_{\pi(i)}\) over all permutations
  \(\pi\); moreover \(S_{\min}<S_{\max}\). The \textbf{average} over all
  permutations equals

  \[\frac1{n!}\sum_{\pi}\sum_i x'_i y'_{\pi(i)}  
    = \sum_i x'_i \Big(\frac1{n!}\sum_{\pi} y'_{\pi(i)}\Big)  
    = \sum_i x'_i \cdot \frac1n\sum_j y'_j  
    = 0.\]

  Since the minimum is \(<\) the maximum and the average is \(0\), it
  follows that \(S_{\min}<0<S_{\max}\). Hence \(r_{\min}<0<r_{\max}\).

  Conversely, if \(\sigma_X=0\) or \(\sigma_Y=0\), then one centered
  sequence is identically zero, so \(S_{\max}=S_{\min}=0\) and the
  correlations are undefined or (in a limiting sense) attain 0. This
  matches the ``only if'' direction.
\end{itemize}

\begin{enumerate}
\def\labelenumi{\Alph{enumi})}
\setcounter{enumi}{1}
\tightlist
\item
  Population (FH/quantile) setting
\end{enumerate}

Let \(Q_X, Q_Y : [0,1]\to\mathbb{R}\) be the quantile functions of the
\textbf{scored} variables. Set

\[f(u):=Q_X(u)-\mu_X,\qquad g(u):=Q_Y(u)-\mu_Y,\]

so \(\int_0^1 f=\int_0^1 g=0\). Monotone scoring preserves category
order, hence \(Q_X, Q_Y\) are nondecreasing; therefore \(f, g\) are
nondecreasing functions.

\begin{itemize}
\item
  \textbf{Upper/comonotone covariance}:

  \[\mathrm{Cov}_{\max}=\int_0^1 f(u)\,g(u)\,du \;\ge\; \Big(\int_0^1 f\Big)\Big(\int_0^1 g\Big)=0\]

  by Chebyshev's \textbf{integral} inequality for similarly ordered
  functions.
\item
  \textbf{Lower/antitone covariance}:

  \[\mathrm{Cov}_{\min}=\int_0^1 f(u)\,g(1-u)\,du \;\le\; \Big(\int_0^1 f\Big)\Big(\int_0^1 g(1-u)\Big)=0,\]

  since \(f\) is nondecreasing and \(g(1-u)\) is \textbf{nonincreasing}.
\end{itemize}

Dividing by \(\sigma_X\sigma_Y>0\) yields \(r_{\max}\ge 0\) and
\(r_{\min}\le 0\).

\begin{itemize}
\tightlist
\item
  \textbf{Strictness.}\\
  If both variables have positive variance, then \(f\) and \(g\) are not
  a.e. constant. For nontrivial nondecreasing \(f,g\), Chebyshev's
  integral inequalities above are \textbf{strict} unless one function is
  a.e. constant. Hence \(\mathrm{Cov}_{\max}>0\) and
  \(\mathrm{Cov}_{\min}<0\), giving \(r_{\max}>0\) and \(r_{\min}<0\).
  Conversely, if either variance is \(0\), both covariances are \(0\).
\end{itemize}

This completes the proof. \(\square\)
\end{proof}

\begin{center}\rule{0.5\linewidth}{0.5pt}\end{center}

\subsubsection{remarks (orientation \&
invariance)}\label{remarks-orientation-invariance}

\begin{itemize}
\item
  The lemma depends only on \textbf{order}, not spacing. Any
  (non-decreasing) rescoring---including affine changes like
  \(10,20,30,\dots\) or nonlinear monotone spacings---preserves the sign
  conclusions.
\item
  If you encode one variable with a \textbf{decreasing} scoring map,
  then ``comonotone'' in category order becomes \textbf{antitone} in
  numeric order; swap the labels to recover the same sign statements.
\item
  For \textbf{Spearman}, take scores as midranks; everything above
  applies verbatim.
\end{itemize}

\subsection{5.4 Symmetry condition}\label{symmetry-condition}

\begin{itemize}
\tightlist
\item
  If one marginal is \textbf{palindromically symmetric} about its center
  rank, then \(r_{\min} = -r_{\max}\).\\
  \emph{(Cees -- r\_min MAGNITUDE Asymmetry in Pearson Correlation.md)}
\end{itemize}

\begin{center}\rule{0.5\linewidth}{0.5pt}\end{center}

\subsection{necessary/sufficient symmetry condition (precisely
stated)}\label{necessarysufficient-symmetry-condition-precisely-stated}

\begin{itemize}
\item
  \textbf{Pair-specific equality} \(|r_{\min}|=r_{\max}\) \(\iff\)
  \(\sum_i x'_i\,(y'_i+y'_{n+1-i})=0\).\\
  (Can happen by coincidence even if \(Y\) isn't symmetric.)
\item
  \textbf{Universal in \(X\)} equality (for \textbf{all} \(X\)):\\
  \textbf{Necessary and sufficient} that \textbf{one} marginal (say
  \(Y\)) has

  \begin{enumerate}
  \def\labelenumi{(\roman{enumi})}
  \tightlist
  \item
    \textbf{palindromic mass symmetry} \(p_Y(k)=p_Y(K_Y-1-k)\) and\\
  \item
    \textbf{antisymmetric scoring} \(s_Y(k)+s_Y(K_Y-1-k)=\text{const}\)
    (true for any affine map of \(0{:}(K_Y-1)\)).\\
    This condition is \textbf{independent of \(K_X\) vs \(K_Y\)}.
  \end{enumerate}
\end{itemize}

\subsection{Symmetry}\label{symmetry}

\subsubsection{1) setup (permutation extremes for a single
sample)}\label{setup-permutation-extremes-for-a-single-sample}

Let \(x_{(1)}\ge\cdots\ge x_{(n)}\) and \(y_{(1)}\ge\cdots\ge y_{(n)}\)
be the sorted sample values (or expanded from counts). Define centered
sequences

\[x'_i:=x_{(i)}-\bar x,\qquad y'_i:=y_{(i)}-\bar y,\]

and the two permutation extremes

\[S_{\max}=\sum_{i=1}^n x'_i\,y'_i,\qquad  
S_{\min}=\sum_{i=1}^n x'_i\,y'_{n+1-i},\]

with \(r_{\max}=S_{\max}/(\sigma_X\sigma_Y)\),
\(r_{\min}=S_{\min}/(\sigma_X\sigma_Y)\).

We always have \(S_{\max}\ge 0\) and \(S_{\min}\le 0\). Magnitude
symmetry \(|r_{\min}|=r_{\max}\) is equivalent to

\[S_{\max}+S_{\min}=0.\]

\subsubsection{\texorpdfstring{2) \textbf{pair-specific necessary \&
sufficient
condition}}{2) pair-specific necessary \& sufficient condition}}\label{pair-specific-necessary-sufficient-condition}

Let

\[z_i := y'_i + y'_{n+1-i}.\]

Then

\[S_{\max}+S_{\min}=\sum_{i=1}^n x'_i\,z_i.\]

Hence, for this \emph{particular} pair \((x,y)\),

\[\boxed{\ |r_{\min}|=r_{\max}\ \Longleftrightarrow\ \sum_i x'_i\,z_i=0\ } \tag{★}\]

(i.e., the centered \(X\) sequence is orthogonal to the
``palindrome-sum'' vector \(z\)).\\
This can hold even if \(Y\) is not symmetric---by coincidence of that
inner product being zero.

\subsubsection{\texorpdfstring{3) \textbf{universal condition
(independent of
\(X\))}}{3) universal condition (independent of X)}}\label{universal-condition-independent-of-x}

If you want \(|r_{\min}|=r_{\max}\) to hold for \textbf{every} possible
\(X\) (every centered, sorted \(x'\)), then (★) must hold for all
\(x'\). The only way is

\[z_i \equiv 0\quad\text{for all }i,\]

i.e.

\[y'_i + y'_{n+1-i}\equiv 0\ \ \ \forall i.\]

Translating back to categories and scores, a \textbf{necessary and
sufficient universal condition} is:

\begin{itemize}
\item
  \textbf{Palindromic mass symmetry for \(Y\):} \(p_Y(k)=p_Y(K_Y-1-k)\)
  for all \(k\), and
\item
  \textbf{Antisymmetric scoring for \(Y\):} the score map satisfies
  \(s_Y(k)+s_Y(K_Y-1-k)\) is constant (true for affine scores like
  \(0,1,\dots,K_Y-1\), and for Spearman midranks under palindromic
  counts).
\end{itemize}

Under these, \(z_i\equiv 0\) and thus \(|r_{\min}|=r_{\max}\) for any
\(X\).

\begin{quote}
If either mass symmetry or score antisymmetry fails, you can generally
find an \(X\) for which \(|r_{\min}|\ne r_{\max}\) (often
\(|r_{\min}|>r_{\max}\) when \(Y\) is skewed).
\end{quote}

\subsubsection{\texorpdfstring{4) does the \textbf{number of categories}
matter?}{4) does the number of categories matter?}}\label{does-the-number-of-categories-matter}

\begin{itemize}
\item
  \textbf{For equality itself:} No.~Odd vs even, and equal vs unequal
  numbers of categories, do \textbf{not} affect whether
  \(|r_{\min}|=r_{\max}\) can hold. What matters is the
  symmetry/antisymmetry described above.
\item
  \textbf{For the sizes of the extremes:} Different \(K\), skewness, and
  spacing of scores \emph{do} affect how large \(r_{\max}\) and
  \(|r_{\min}|\) can be, but not whether the magnitudes are equal.
\end{itemize}

\subsubsection{5) FH (transport) version}\label{fh-transport-version}

For FH extremes (optimize over all joints with given marginals), replace
sums by integrals over quantiles:

\[\mathrm{Cov}_{\max}=\int_0^1 (Q_X(u)-\mu_X)(Q_Y(u)-\mu_Y)\,du,\quad  
\mathrm{Cov}_{\min}=\int_0^1 (Q_X(u)-\mu_X)(Q_Y(1-u)-\mu_Y)\,du.\]

Then

\[\mathrm{Cov}_{\max}+\mathrm{Cov}_{\min}=\int_0^1 (Q_X(u)-\mu_X)\,\bigl[Q_Y(u)+Q_Y(1-u)-2\mu_Y\bigr]\,du.\]

So the \textbf{pair-specific} condition is
\(\int (Q_X-\mu_X)\,[Q_Y(u)+Q_Y(1-u)-2\mu_Y]\,du=0\).\\
The \textbf{universal (in \(X\))} condition is \(Q_Y(u)+Q_Y(1-u)\)
constant a.e. (palindromic symmetry in distribution + antisymmetric
scoring), exactly mirroring the permutation case.

\subsubsection{6) practical diagnostic}\label{practical-diagnostic}

Given counts and a fixed score map:

\begin{itemize}
\item
  Expand to sequences (or evaluate on a fine quantile grid for FH),
\item
  Form \(x'\), \(y'\), and \(z=y'+\text{rev}(y')\),
\item
  Check \(\sum x' z\).

  \begin{itemize}
  \item
    If \(\approx 0\) ⇒ expect \(|r_{\min}|\approx r_{\max}\).
  \item
    If \(<0\) ⇒ often \(|r_{\min}|>r_{\max}\) (your BES pattern).
  \item
    If \(>0\) ⇒ often \(|r_{\min}|<r_{\max}\).
  \end{itemize}
\end{itemize}

\begin{center}\rule{0.5\linewidth}{0.5pt}\end{center}

\subsubsection{tl;dr}\label{tldr}

\begin{itemize}
\item
  \textbf{Pair-specific necessity \& sufficiency:}
  \(\sum_i x'_i\,(y'_i+y'_{n+1-i})=0\).
\item
  \textbf{Universal (in \(X\)) necessity \& sufficiency:} one marginal
  (say \(Y\)) has \textbf{palindromic mass symmetry} and
  \textbf{antisymmetric scores}---then \(|r_{\min}|=r_{\max}\) for all
  \(X\).
\end{itemize}

\subsubsection{Examples}\label{examples}

Below, ``scores \(0{:}(K-1)\)'' means category values
\([0,1,\dots,K-1]\). ``Affine recode'' means
\(a+b\times\text{(old score)}\) with \(b>0\). ``Non-affine monotone''
means strictly increasing but not affine (e.g., convex spacing).

\subsubsection{Notation (permutation extremes, ``Lane
A'')}\label{notation-permutation-extremes-lane-a}

\begin{itemize}
\item
  Build expanded vectors \(x\) and \(y\) from \textbf{counts} and
  \textbf{scores}.
\item
  \(r_{\max}=\mathrm{cor}(\text{sort}(x,\downarrow),\text{sort}(y,\downarrow))\).
\item
  \(r_{\min}=\mathrm{cor}(\text{sort}(x,\downarrow),\text{sort}(y,\uparrow))\).
\end{itemize}

\begin{center}\rule{0.5\linewidth}{0.5pt}\end{center}

\subsubsection{\texorpdfstring{A) \textbf{Same number of categories, but
\(|r_{\min}|\ne r_{\max}\)} (concrete
numbers)}{A) Same number of categories, but \textbar r\_\{\textbackslash min\}\textbar\textbackslash ne r\_\{\textbackslash max\} (concrete numbers)}}\label{a-same-number-of-categories-but-r_minne-r_max-concrete-numbers}

\textbf{Counts (both \[K=5\]):}

\begin{itemize}
\item
  \[X:\ [100,\,200,\,300,\,200,\,200]\]
\item
  \(Y:\ [100,\,100,\,100,\,300,\,400]\) (skewed; not palindromic)
\end{itemize}

\textbf{Scores:} \(0{:}4\) for both.

\begin{itemize}
\item
  \[r_{\max}\approx 0.929398758\]
\item
  \[r_{\min}\approx -0.881118303\]
\item
  \[|r_{\min}| \approx 0.8811 \ne r_{\max}\]
\end{itemize}

\textbf{Affine recode (10,20,30,40,50) for both:}

\begin{itemize}
\tightlist
\item
  \(r_{\max}, r_{\min}\) \textbf{unchanged} (up to rounding).
\end{itemize}

\textbf{Non-affine monotone recode for \(Y\):} \([0,1,1.5,3,10]\).

\begin{itemize}
\item
  \[r_{\max}\approx 0.930261031\]
\item
  \[r_{\min}\approx -0.863250703\]
\item
  Values change (Pearson is \textbf{not} invariant to non-affine
  spacing); magnitudes still unequal.
\end{itemize}

🧠 \textbf{Why unequal magnitudes?} No palindromic symmetry in \(Y\) ⇒
the ``palindrome-sum'' pattern in \(y'-\)space isn't zero, so
\(S_{\max}+S_{\min}\ne 0\).

\begin{center}\rule{0.5\linewidth}{0.5pt}\end{center}

\subsubsection{\texorpdfstring{B) \textbf{Different numbers of
categories, yet \(|r_{\min}|=r_{\max}\)} (concrete
numbers)}{B) Different numbers of categories, yet \textbar r\_\{\textbackslash min\}\textbar=r\_\{\textbackslash max\} (concrete numbers)}}\label{b-different-numbers-of-categories-yet-r_minr_max-concrete-numbers}

\textbf{Counts:}

\begin{itemize}
\item
  \(X\) (4 categories): \([100,\,200,\,300,\,400]\)
\item
  \(Y\) (5 categories): \([100,\,200,\,400,\,200,\,100]\)
  (\textbf{palindromic})
\end{itemize}

\textbf{Scores:} \(X:\ 0{:}3\), \(Y:\ 0{:}4\).

\begin{itemize}
\item
  \[r_{\max}\approx 0.912870929\]
\item
  \[r_{\min}\approx -0.912870929\]
\item
  \(|r_{\min}| = r_{\max}\) (to machine precision)
\end{itemize}

\textbf{Affine recode (e.g., 10,20,\ldots):}

\begin{itemize}
\tightlist
\item
  Magnitudes remain \textbf{identical}; affine transforms don't change
  Pearson correlation.
\end{itemize}

\textbf{Non-affine monotone recode for \(Y\):} \([0,1,1.5,3,10]\).

\begin{itemize}
\item
  \[r_{\max}\approx 0.580072921\]
\item
  \[r_{\min}\approx -0.842041337\]
\item
  \textbf{Magnitudes no longer match} because antisymmetry of the
  \textbf{scores} about the center is lost, even though the
  \textbf{counts} are palindromic.
\end{itemize}

🧠 \textbf{Why equal magnitudes in the affine case?} Palindromic
\textbf{mass} symmetry in \(Y\) + \textbf{antisymmetric scoring} (any
affine map of \(0{:}(K_Y-1)\)) ⇒ the palindrome-sum vector
\(y'_i+y'_{n+1-i}\equiv 0\). Hence \(S_{\max}+S_{\min}=0\) ⇒
\(|r_{\min}|=r_{\max}\) for \textbf{any} \(X\), irrespective of
\(K_X\neq K_Y\).

\begin{center}\rule{0.5\linewidth}{0.5pt}\end{center}

\subsubsection{\texorpdfstring{C) \textbf{Spearman invariance (same data
as in
B)}}{C) Spearman invariance (same data as in B)}}\label{c-spearman-invariance-same-data-as-in-b}

Compute ranks (with average ties) \textbf{once} for \(x\) and \(y\),
then form extremes by sorting those rank vectors.

\begin{itemize}
\item
  With \(Y\) scored \(0{:}4\):\\
  \(r_{\max}^{\text{Spearman}}\approx 0.920837215,\ \ r_{\min}^{\text{Spearman}}\approx -0.920837215\).
\item
  With \(Y\) scored non-affinely \([0,1,1.5,3,10]\):\\
  \textbf{identical} Spearman extremes: \(0.920837215\) and
  \(-0.920837215\).
\end{itemize}

✅ \textbf{Spearman is invariant} to any strictly monotone recode
because ranks depend only on order, not spacing.

\begin{center}\rule{0.5\linewidth}{0.5pt}\end{center}

\subsection{what changes and what doesn't (Pearson vs
Spearman)}\label{what-changes-and-what-doesnt-pearson-vs-spearman}

\begin{itemize}
\item
  \textbf{Pearson:}

  \begin{itemize}
  \item
    \textbf{Affine} recodes \(a+b\cdot\text{score}\) (with \(b>0\))
    leave \(r_{\max}, r_{\min}\), and any \(|r_{\min}|=r_{\max}\)
    conclusions \textbf{unchanged}.
  \item
    \textbf{Non-affine monotone} recodes can change
    \(r_{\max}, r_{\min}\) and can break \(|r_{\min}|=r_{\max}\) even
    with palindromic counts.
  \end{itemize}
\item
  \textbf{Spearman:}

  \begin{itemize}
  \item
    Any strictly monotone recode leaves extremes \textbf{unchanged}.
  \item
    Palindromic counts in \(Y\) (with the usual symmetric rank scheme)
    give \(|r_{\min}|=r_{\max}\).
  \end{itemize}
\end{itemize}

\begin{center}\rule{0.5\linewidth}{0.5pt}\end{center}

\subsection{6. Simulation /
Randomization}\label{simulation-randomization}

\textbf{Permutation distribution:}\\
Generate full vectors from marginals, independently permute \(X\) and
\(Y\), compute \(r\), repeat. Produces null distribution centered at
0.\\
\emph{(Proofs-OrdinalVariables.md)}

\begin{center}\rule{0.5\linewidth}{0.5pt}\end{center}

\subsection{7. Key Tools Referenced}\label{key-tools-referenced}

\begin{itemize}
\tightlist
\item
  \textbf{Fréchet--Hoeffding bounds} for extremal joint CDFs.\\
  \emph{(BFbounds.md)}
\item
  \textbf{Hardy--Littlewood--Pólya rearrangement inequality} for
  rank-based maximization/minimization.\\
  \emph{(Ordinal Variables and Correlation.md)}
\item
  \textbf{Comonotone/countermonotone quantile coupling}.\\
  \emph{(Math Proof GPT -- Whitt 1976 Paper.md)}
\item
  \textbf{Capacity limits} in discrete joint tables explain why FH lower
  may be unattainable.
\end{itemize}

\begin{center}\rule{0.5\linewidth}{0.5pt}\end{center}

\subsection{8. Master Takeaways}\label{master-takeaways}

\begin{enumerate}
\def\labelenumi{\arabic{enumi}.}
\tightlist
\item
  \textbf{\(r_{\max}\) = FH upper (comonotonic), always attainable} in
  discrete ordinal settings.
\item
  \textbf{\(r_{\min}\)} may \textbf{not} equal FH lower in discrete
  settings; use best feasible opposite-order arrangement.
\item
  No universal bound on relative magnitudes: \(|r_{\min}|\) may be less
  than, equal to, or greater than \(r_{\max}\).
\item
  Palindromic symmetry of one marginal guarantees
  \(|r_{\min}| = r_{\max}|\).
\item
  Binary cases illustrate when perfect alignment is feasible and when
  skew prevents it.
\end{enumerate}

\begin{center}\rule{0.5\linewidth}{0.5pt}\end{center}

\newpage

\section{Positive Semi-Definiteness in Covariance Matrices: Implications
for PCA, Factor Analysis, and
SEM}\label{positive-semi-definiteness-in-covariance-matrices-implications-for-pca-factor-analysis-and-sem}

\subsection{1. Variance--Covariance Matrices: Definition and Core
Property}\label{variancecovariance-matrices-definition-and-core-property}

\paragraph{1.1 Definition}\label{definition}

For a random vector \(\mathbf{X} = (X_1, X_2, \dots, X_n)^\top\), the
variance--covariance matrix is

\[
\mathbf{\Sigma} = \mathbb{E} \left[ (\mathbf{X} - \mathbb{E}[\mathbf{X}]) (\mathbf{X} - \mathbb{E}[\mathbf{X}])^\top \right],
\] where \(\sigma_{ij} = \operatorname{Cov}(X_i, X_j)\).\\
\emph{(vcov\_PSD.md):contentReference{oaicite:2}}

\begin{itemize}
\tightlist
\item
  \textbf{Diagonal elements} \(\sigma_{ii}\) are variances, always
  \(\ge 0\).
\item
  \textbf{Off-diagonal} entries are covariances.
\end{itemize}

\begin{center}\rule{0.5\linewidth}{0.5pt}\end{center}

\paragraph{1.2 PSD Property}\label{psd-property}

A matrix is \textbf{positive semi-definite (PSD)} if, for all nonzero
\(\mathbf{z}\), \[
\mathbf{z}^\top \mathbf{\Sigma} \mathbf{z} \ge 0.
\] For covariance matrices: \[
\mathbf{z}^\top \mathbf{\Sigma} \mathbf{z}
= \mathbb{E}\!\left[(\mathbf{z}^\top (\mathbf{X} - \mathbb{E}[\mathbf{X}]))^2\right] \ge 0,
\] since it is an expectation of a square.\\
\emph{(vcov\_PSD.md):contentReference{oaicite:3}}

\begin{center}\rule{0.5\linewidth}{0.5pt}\end{center}

\paragraph{1.3 Positive Definite (PD) vs
PSD}\label{positive-definite-pd-vs-psd}

\begin{itemize}
\tightlist
\item
  \textbf{PD} if \(\mathbf{z}^\top \mathbf{\Sigma} \mathbf{z} > 0\) for
  all \(\mathbf{z} \neq 0\), i.e., no perfect linear dependence among
  variables.
\item
  \textbf{PSD} if at least one nonzero \(\mathbf{z}\) yields equality
  (rank-deficient case).
  \emph{(vcov\_PSD.md):contentReference{oaicite:4}}
\end{itemize}

\begin{center}\rule{0.5\linewidth}{0.5pt}\end{center}

\subsection{2. PCA and PSD Requirements}\label{pca-and-psd-requirements}

\paragraph{2.1 PCA Overview}\label{pca-overview}

Principal Components Analysis (PCA) decomposes a sample covariance or
correlation matrix \(\mathbf{S} = n^{-1} \mathbf{X}^\top \mathbf{X}\)
into \[
\mathbf{S} = \mathbf{Q} \,\mathrm{diag}(\lambda_1,\dots,\lambda_p)\, \mathbf{Q}^\top.
\] \emph{(PCA and FA PSD requirement.md):contentReference{oaicite:5}}

\begin{center}\rule{0.5\linewidth}{0.5pt}\end{center}

\paragraph{2.2 Why PSD Matters for PCA}\label{why-psd-matters-for-pca}

\begin{itemize}
\tightlist
\item
  \textbf{Mathematically}: Eigen-decomposition works for any symmetric
  \(\mathbf{S}\), PSD or not.
\item
  \textbf{Interpretively}: Negative eigenvalues break the ``variance
  explained'' story (variance cannot be negative).
\item
  \textbf{Software behavior}: Most PCA routines check PD/PSD; failure
  can cause warnings or dropped components. \emph{(PCA and FA PSD
  requirement.md):contentReference{oaicite:6}}
\end{itemize}

\begin{center}\rule{0.5\linewidth}{0.5pt}\end{center}

\paragraph{2.3 If Input is Symmetric but not
PSD}\label{if-input-is-symmetric-but-not-psd}

Example:\\
\[
\mathbf{A} = \begin{bmatrix} 1 & 2 \\ 2 & 1 \end{bmatrix},\quad \lambda = (3,-1),
\] is indefinite (one negative eigenvalue).\\
- PCA math still runs; eigenvalues/vectors computed. - Interpretation
fails: \(\sqrt{\lambda_j}\) imaginary for \(\lambda_j<0\). - Some
software refuses input; others run but return nonsensical variance
metrics. \emph{(PCA and FA PSD
requirement.md):contentReference{oaicite:7}}

\begin{center}\rule{0.5\linewidth}{0.5pt}\end{center}

\paragraph{2.4 Remedies for PCA}\label{remedies-for-pca}

\begin{enumerate}
\def\labelenumi{\arabic{enumi}.}
\tightlist
\item
  Treat tiny negative eigenvalues as zero.\\
\item
  Adjust to nearest PSD matrix (Higham 2002).\\
\item
  Run SVD on the raw data matrix instead of precomputed
  \(\mathbf{S}\).\\
  \emph{(PCA and FA PSD requirement.md):contentReference{oaicite:8}}
\end{enumerate}

\begin{center}\rule{0.5\linewidth}{0.5pt}\end{center}

\subsection{3. Factor Analysis and PSD
Requirements}\label{factor-analysis-and-psd-requirements}

\paragraph{3.1 FA Model}\label{fa-model}

Common factor analysis models \(\boldsymbol\Sigma\) as: \[
\boldsymbol\Sigma = \boldsymbol\Lambda \boldsymbol\Lambda^\top + \boldsymbol\Psi,
\] where \(\boldsymbol\Psi\) is diagonal with positive entries.
\emph{(PCA and FA PSD requirement.md):contentReference{oaicite:9}}

\begin{center}\rule{0.5\linewidth}{0.5pt}\end{center}

\paragraph{3.2 PSD/PD Needs in FA}\label{psdpd-needs-in-fa}

\begin{itemize}
\tightlist
\item
  \textbf{ML FA}: Requires \(\mathbf{S} \succ 0\);
  \(\log \det \mathbf{S}\) undefined if indefinite.
\item
  \textbf{Principal-axis / Image FA}: Also break with negative
  eigenvalues (square roots fail). \emph{(PCA and FA PSD
  requirement.md):contentReference{oaicite:10}}
\end{itemize}

\begin{center}\rule{0.5\linewidth}{0.5pt}\end{center}

\paragraph{3.3 Remedies for FA}\label{remedies-for-fa}

\begin{itemize}
\tightlist
\item
  Repair \(\mathbf{S}\) to PSD before analysis.
\item
  Add small ridge term \(\epsilon \mathbf{I}\). \emph{(PCA and FA PSD
  requirement.md):contentReference{oaicite:11}}
\end{itemize}

\begin{center}\rule{0.5\linewidth}{0.5pt}\end{center}

\subsection{4. SEM and PSD Requirements}\label{sem-and-psd-requirements}

\paragraph{4.1 SEM Machinery}\label{sem-machinery}

Observed covariance \(\mathbf{S}\); implied covariance
\(\boldsymbol\Sigma(\boldsymbol\theta)\) from structural and measurement
models. \emph{(PCA and FA PSD
requirement.md):contentReference{oaicite:12}}

\begin{center}\rule{0.5\linewidth}{0.5pt}\end{center}

\paragraph{4.2 Fit Functions and PD
Needs}\label{fit-functions-and-pd-needs}

\begin{itemize}
\tightlist
\item
  \textbf{ML / GLS}: Require both \(\mathbf{S}\) and
  \(\boldsymbol\Sigma(\theta)\) PD.
\item
  \textbf{WLS/DWLS}: Need PD indirectly (weight matrix invertible).
\item
  \textbf{ULS}: Can start from indefinite \(\mathbf{S}\), but test
  statistics again require PD. \emph{(PCA and FA PSD
  requirement.md):contentReference{oaicite:13}}
\end{itemize}

\begin{center}\rule{0.5\linewidth}{0.5pt}\end{center}

\paragraph{4.3 What Happens if Not PSD}\label{what-happens-if-not-psd}

\begin{itemize}
\tightlist
\item
  Software errors out (``matrix is not positive definite'').
\item
  Log-determinant undefined; inversion fails.
\item
  Even if bypassed, optimisation steps hit complex/NaN values.
  \emph{(PCA and FA PSD requirement.md):contentReference{oaicite:14}}
\end{itemize}

\begin{center}\rule{0.5\linewidth}{0.5pt}\end{center}

\paragraph{4.4 Remedies in SEM}\label{remedies-in-sem}

\begin{itemize}
\tightlist
\item
  Nearest PD smoothing (Higham 2002).
\item
  Ridge adjustment.
\item
  Analyse raw data.
\item
  Check rounding issues. \emph{(PCA and FA PSD
  requirement.md):contentReference{oaicite:15}}
\end{itemize}

\begin{center}\rule{0.5\linewidth}{0.5pt}\end{center}

\subsection{5. Core Lessons Across PCA, FA,
SEM}\label{core-lessons-across-pca-fa-sem}

\begin{enumerate}
\def\labelenumi{\arabic{enumi}.}
\tightlist
\item
  \textbf{Theoretical PSD}: Covariance matrices from real data are PSD
  by definition:contentReference{oaicite:16}.
\item
  \textbf{Numerical Issues}: Sampling noise, rounding, or listwise
  deletion can yield small negative
  eigenvalues:contentReference{oaicite:17}.
\item
  \textbf{PCA}: Can compute from indefinite matrix but interpretation
  breaks:contentReference{oaicite:18}.
\item
  \textbf{FA and SEM}: Require PD input; will not run on indefinite
  matrices:contentReference{oaicite:19}.
\item
  \textbf{Repairs}: Adjust to nearest PSD, ridge, or recompute from raw
  data:contentReference{oaicite:20}.
\end{enumerate}

\begin{center}\rule{0.5\linewidth}{0.5pt}\end{center}




\end{document}
